\section{Problema de Pesquisa}
\label{sec:1.2}

A montagem manual de grades horárias em IES é um processo iterativo, sujeito a erros e de elevada carga operacional. Coordenadores precisam conciliar disponibilidade docente, distribuição da carga horária, restrições de turno por curso e ausência de conflitos entre turmas que compartilham professores. Conforme Bashab et al. (2023), à medida que o número de professores, disciplinas, turmas e horários cresce, o espaço de soluções se expande de forma combinatória, tornando a busca manual por uma grade sem conflitos uma tarefa demorada e limitada, sobretudo quando também se busca atender restrições flexíveis.

A classificação do UCTP como problema NP-difícil indica que sua resolução envolve um espaço de busca de crescimento combinatório, no qual o número de soluções possíveis aumenta rapidamente conforme são adicionados professores, turmas, disciplinas, horários e restrições. Em termos práticos, isso significa que não se conhece um algoritmo de tempo polinomial capaz de garantir a solução ótima para todas as instâncias gerais do problema, salvo se $P = NP$ \cite{premananda2024}. Assim, embora métodos exatos possam ser relevantes em cenários menores ou mais controlados, sua aplicação tende a perder viabilidade à medida que aumentam as variáveis envolvidas e as restrições consideradas. Nesse contexto, heurísticas e meta-heurísticas, como Algoritmos Genéticos, Busca Tabu e Simulated Annealing, são amplamente utilizadas para encontrar soluções satisfatórias em tempo computacional aceitável, ainda que sem garantia de ótimo global \cite{eiben2015}.


Embora existam estudos que aplicam meta-heurísticas ao UCTP \cite{ceschia2023, awad2022, elen2022}, muitas soluções permanecem como protótipos acadêmicos isolados, sem integração com sistemas voltados ao uso operacional por coordenadores. Além disso, são menos frequentes comparações entre implementações do mesmo algoritmo em paradigmas de programação distintos dentro de uma arquitetura controlada, o que limita a análise sobre o impacto dessas escolhas no desempenho do sistema.

Diante desse contexto, este trabalho busca responder à seguinte questão de pesquisa: \textit{é possível desenvolver um sistema modular de geração automatizada de grades horárias baseado em Algoritmos Genéticos, integrando múltiplos paradigmas de programação em uma mesma arquitetura, e avaliar comparativamente o desempenho de cada implementação sob condições controladas de benchmark?}
